\documentclass[11pt]{article}
\usepackage[margin=1in]{geometry}
\usepackage{booktabs}
\usepackage{hyperref}

\title{CAI6734 Weekly Report 6: Effect of the Test-Cost Penalty on Training}
\author{Omer Kahveci}
\date{November 3, 2026}

\begin{document}
\maketitle

\section{Project Objective}

The Week 5 report documented the training and evaluation of the
Meta-Llama-3-8B-Instruct LA-CDM model. This week's objective was to examine
how changes to the training configuration affect the improvements obtained
from training, beginning with the test-cost penalty in the training reward.
This penalty discourages the Decision Agent from ordering diagnostic tests;
adjusting its weight is therefore expected to shift the balance between
diagnostic accuracy and diagnostic cost. The initial hypothesis was that a
stronger penalty would reduce cost while largely preserving accuracy. Each
variant is compared with the Week 5 trained model (68.75\% test accuracy and
\$1,137.67 mean diagnostic cost) on the same 240-case test split, using one
rollout per case.

\section{Accomplishments}

\subsection*{Experimental Setup}

A multiplier was added to the test-cost term of the training reward. The
multiplier scales the per-test penalties and leaves the diagnosis and
action-format rewards unchanged; the Week 5 model corresponds to a multiplier
of 1.0. Three variants were planned: 0x (no penalty), 1.5x, and 2x. Each
variant uses the same base model, random seed, data splits, 300-step training
schedule, and remaining settings as the Week 5 run, so that the multiplier is
the only experimental change. The dollar costs used to report diagnostic cost
during evaluation are identical across variants. The 0x and 1.5x variants were
trained this week.

\subsection*{Zero-Cost Variant}

The 0x variant completed training and was evaluated on the test split.
Table~\ref{tab:test} compares its results with those of the Week 5 trained
model.

\begin{table}[ht]
  \centering
  \caption{Test results on the same 240 cases, with one rollout per case.}
  \label{tab:test}
  \begin{tabular}{lrr}
    \toprule
    Metric & Original (1.0x) & Zero-cost (0x) \\
    \midrule
    Diagnosis accuracy, all predictions & 68.75\% & 77.50\% \\
    Diagnosis accuracy, diagnosed only & 70.21\% & 81.22\% \\
    Undiagnosed fraction & 2.08\% & 4.58\% \\
    Action-format reward & 0.9792 & 0.9542 \\
    Macro F1, all predictions & 66.27\% & 74.50\% \\
    Hypothesis accuracy & 63.26\% & 72.97\% \\
    Expected calibration error & 10.38\% & 8.77\% \\
    Average tests requested & 2.73 & 2.29 \\
    Average diagnostic cost & \$1,137.67 & \$1,627.71 \\
    \bottomrule
  \end{tabular}
\end{table}

Removing the penalty increased overall test accuracy by 8.75 percentage points
(186 versus 165 correct diagnoses out of 240). Macro F1, hypothesis accuracy,
and expected calibration error also improved. However, mean diagnostic cost
increased by \$490.04 (approximately 43\%), even though the model requested
approximately 0.44 fewer tests per case. This pattern suggests a shift toward
more expensive tests, although the aggregate metrics do not identify which
tests changed. The number of undiagnosed cases increased from 5 to 11. Because
these results come from a single training run, the magnitude of the
improvement has not yet been established across random seeds.

\subsection*{1.5x Variant}

The 1.5x variant completed training, but the resulting model did not perform
as intended, and the cause is under investigation. In the final validation
pass, none of the 960 rollouts (240 cases with four rollouts each) produced a
valid diagnosis, and the action-format reward was 0. The model was therefore
not evaluated on the test split. To determine when the problem began, a script
was written to summarize the training rewards in each of the three training
phases (Table~\ref{tab:collapse}).

\begin{table}[ht]
  \centering
  \caption{1.5x variant: training steps with a positive batch-mean reward.}
  \label{tab:collapse}
  \begin{tabular}{lrrr}
    \toprule
    Training phase & Format & Diagnosis & Mean cost reward \\
    \midrule
    Decision Agent GRPO, steps 1--100 & 1/100 & 1/100 & $-0.0553$ \\
    Hypothesis SFT, steps 101--200 & 0/100 & 0/100 & $-0.0259$ \\
    Confidence calibration, steps 201--300 & 0/100 & 0/100 & $-0.0213$ \\
    \bottomrule
  \end{tabular}
\end{table}

The failure was already present in the initial Decision Agent training steps,
before the Hypothesis Agent phases began. The transition between training
phases therefore does not account for it.

\section{Challenges}

The principal open issue is the failure of the 1.5x variant. One possible
explanation is that the untrained Llama model rarely produces a valid action
at the start of training, so the diagnosis and format rewards are almost
always zero during the early steps. Because GRPO compares rewards within each
group of four rollouts, the early learning signal would then be dominated by
differences in test cost. This could push the policy toward ordering fewer
tests rather than toward producing valid diagnoses. This explanation has not
been confirmed; the small group size and the model's initial format errors may
also contribute.

In addition, the 0x result ran counter to the initial hypothesis: removing the
penalty improved accuracy, at the expense of higher diagnostic cost. The logs
for both variants also contain repeated requests for an unrecognized
\texttt{Complete Blood Count (CBC)} action, together with invalid-format
warnings, which have not yet been traced.

\section{Next Steps}

\begin{itemize}
  \item Investigate the cause of the 1.5x failure and rerun that variant.
  \item Update the evaluation pipeline to record every decision the model makes
        on the validation set in a CSV file, with one row per action. This will
        show the number of tests ordered for each patient and facilitate
        inspection of the model's reasoning.
  \item Report 95\% confidence intervals and $p$-values for future results.
        The per-case records will support paired comparisons between models
        evaluated on the same test cases.
  \item Review related literature for approaches to revising the prompts or
        the structure of the agentic loop.
\end{itemize}

\section{AI Disclaimer}

I used OpenAI Codex and Anthropic Claude to assist with implementing the
experiment, inspecting training and evaluation logs, and drafting this report.
I reviewed the reported metrics and remain responsible for the final
interpretation.

\end{document}
